\section{Introduction}\label{sec:intro}

\paragraph{The problem.} Long-term observation of infant posture and movement can support early developmental risk assessment. Automating it starts with 2D keypoint estimation from ordinary RGB images~\citep{huang2021fidip}.
Infants make this a small-data problem in two ways. Privacy limits the number of labelled real images. And infants differ from the adults in COCO~\citep{lin2014coco} in limb proportions, pose repertoire and context, so adult-trained estimators degrade on infant-specific poses~\citep{huang2021fidip}.

\paragraph{FiDIP.} FiDIP (Fine-tuned Domain-adapted Infant Pose)~\citep{huang2021fidip} addresses both. It fine-tunes a COCO-pretrained pose network on SyRIP~\citep{syripdata}, a mix of 200 real web images and 1000 synthetic renders, while a domain classifier pushes the features of real and synthetic images together.
Its strongest variant, DarkPose~\citep{zhang2020darkpose} with an HRNet-W48 backbone~\citep{sun2019hrnet}, reports 93.6~AP on SyRIP Test100~\citep{fidipcode}. As far as we could establish (\cref{sec:related}), this is still the best published number under that protocol.

\paragraph{Goals.} This report documents a hands-on reproduction of FiDIP and an attempt to improve on it, carried out with an AI coding agent on Northeastern's Explorer cluster. We had four goals:
\begin{enumerate}[label=\textbf{(G\arabic*)}, leftmargin=*]
  \item Reproduce the published baseline, from the released checkpoints and by retraining.
  \item Make the evaluation conditions explicit, and identical for every number we compare.
  \item Improve Test100 AP on the same split, without using Test100 for any decision.
  \item Keep the work reproducible and the agentic process auditable.
\end{enumerate}

\paragraph{Findings.}
\begin{itemize}[leftmargin=*]
  \item \textbf{Reproduction} (\cref{sec:repro}). The released FiDIP checkpoint scores \res{rel.fidip.test.ap}~AP on Test100 with our evaluator, against the published 93.6.
  Retraining with the released recipe gives \resms{rt.fidip.test.ap}~AP with domain adaptation and \resms{rt.ft.test.ap}~AP without it (\res{rt.seeds} seeds each).
  \item \textbf{Protocol} (\cref{sec:fidip,sec:data}). Run literally, the README training command fine-tunes from random weights with the adaptation weight $\lambda=0$, and the training loop keeps the epoch that is best on Test100.
  We select every model on Val400, the 397 validate500 images that neither belong to nor duplicate Test100, and pre-register each configuration before its single Test100 evaluation.
  \item \textbf{Data and evaluation conditions} (\cref{sec:data,sec:eval}). Hashing, near-duplicate search and label/image consistency checks found six issues that would silently bias results:
  \begin{itemize}[leftmargin=*]
    \item a constant placeholder object area in validate500;
    \item an annotation with id~0 in Test100 that pycocotools never matches, which caps AR at 0.99 and moves AP by \res{sens.id0.delta} for the released model;
    \item five Test100 images that duplicate training images;
    \item three Val400 candidates that duplicate Test100 images;
    \item MINI-RGBD labels on the synthetic images in the README's training file;
    \item unannotated real images that duplicate validation images.
  \end{itemize}
  MMPose's evaluator lacks the id-0 artifact, so we accept a claim only if it holds under both matching protocols.
  \item \textbf{Improvement} (\cref{sec:improve}). We ran a controlled backbone and input-resolution comparison under one MMPose recipe, plus ablations, all selected on Val400.
  The pre-registered winner, \res{best.name}, scores \resms{best.test.ap}~AP on Test100. That is \res{best.delta.ap}~AP from the released checkpoint, with a paired-bootstrap 95\% CI of [\res{best.delta.ci.lo}, \res{best.delta.ci.hi}].
  \item \textbf{Process} (\cref{sec:lessons,sec:aiuse}). The agent kept \code{SKILL.md}, \code{AGENT.md} and \code{MANIFEST.md} current during the work, not after it, and every reported number traces to a run and a commit.
\end{itemize}

\paragraph{Roadmap.}
\begin{itemize}[leftmargin=*]
  \item \Cref{sec:background}: FiDIP and SyRIP, including where the released code differs from the paper.
  \item \Cref{sec:impl}: the cluster workflow, data preparation, split protocol and evaluator.
  \item \Cref{sec:repro,sec:improve,sec:analysis}: reproduction, improvement and analysis.
  \item \Cref{sec:limits,sec:lessons,sec:aiuse}: limitations, lessons learned and AI-tool use.
  \item The appendix: commands, hyperparameters, compute, provenance and a rubric map.
\end{itemize}
